% SPDX-License-Identifier: CC-BY-NC-ND-4.0
% Copyright (c) 2026 Ingolf Lohmann.
% Current arXiv successor candidate.  Local staging only; not submitted.
\documentclass[11pt]{article}

\usepackage[T1]{fontenc}
\usepackage[utf8]{inputenc}
\usepackage{lmodern}
\usepackage{microtype}
\usepackage[a4paper,margin=27mm]{geometry}
\usepackage{amsmath,amssymb,amsthm,mathtools}
\usepackage{booktabs,longtable,tabularx,array}
\usepackage{enumitem}
\usepackage{xcolor}
\usepackage{hyperref}
\usepackage{xurl}
\usepackage[nameinlink,noabbrev]{cleveref}

% Keep the PDF trailer identifier deterministic for the frozen source archive.
\ifdefined\pdftrailerid
  \pdftrailerid{<6A1B08F33DB040F9A55E9F3D96F7544B>}
\fi

\hypersetup{
  colorlinks=true,
  linkcolor=blue!50!black,
  citecolor=blue!50!black,
  urlcolor=blue!50!black,
  pdftitle={Observer-Relative Retrocausality: Negative Information Direction Along Monotonic Local Change Time},
  pdfauthor={Ingolf Lohmann},
  pdfsubject={Formal existence result, executable witness, physical reference bridge, and scope boundary},
  pdfkeywords={QIK-VRT, observer-relative retrocausality, local change time, Eigenzeit, provenance, distributed systems, delayed choice, quantum eraser}
}

\setlength{\parindent}{0pt}
\setlength{\parskip}{0.45em}
\setlist[itemize]{leftmargin=1.5em,itemsep=0.15em,topsep=0.25em}
\setlist[enumerate]{leftmargin=1.8em,itemsep=0.2em,topsep=0.25em}

\newtheorem{definition}{Definition}[section]
\newtheorem{theorem}{Theorem}[section]
\newtheorem{proposition}{Proposition}[section]
\newtheorem{corollary}{Corollary}[section]
\theoremstyle{remark}
\newtheorem{remark}{Remark}[section]

\newcommand{\host}{\prec_{\!H}}
\newcommand{\receive}{\prec_{\!O}^{R}}
\newcommand{\sourceorder}{\prec^{\theta}}
\newcommand{\qikvrt}{QIK--VRT}
\newcommand{\zenodo}{\url{https://doi.org/10.5281/zenodo.21888130}}

\title{Observer-Relative Retrocausality:\
\large Negative Information Direction Along Monotonic Local Change Time}
\author{Ingolf Lohmann\\
Independent Researcher, QIK--VRT}
\date{12 August 2026\\Current successor candidate, version 2.0}

\begin{document}
\maketitle

\begin{abstract}
This paper gives the current English synthesis of a QIK--VRT definition of
observer-relative retrocausality.  Time is first operationally visible to a
system as an ordered sequence of effective local changes.  QIK--VRT calls a
strictly increasing parameter on that sequence \emph{Eigenzeit}; for a
physical observer on a timelike worldline it may additionally be calibrated
to metric proper time.  Consider two authentic, information-bearing records.
If receiver-local change time increases while the provenance-bound source
marks of the newly received records decrease, then the receiver has negative
information direction:
\[
  \Delta\tau_R>0,\qquad \Delta\theta<0.
\]
Equivalently, the reception order is inverse to the source order for the
record pair.  A finite theorem proves this relation under explicit order,
provenance, information, and monotonicity assumptions.  An executable witness
checks one concrete instance, including strict evidence refinement, a second
observer with the opposite orientation, positive path delays, and a
no-signalling marginal check.

The result does not say that metric proper time runs backwards, that a record
arrives before its own emission, that the past is overwritten, or that a
controllable message can be sent into the past.  A physical apparatus that
satisfies the stated reference bridge instantiates the defined operational
relation.  Delayed-choice and quantum-eraser experiments supply a related
empirical bridge: later context can make a correlation class of earlier
records available without changing earlier raw records or enabling a local
backward channel.  The paper retains Ingolf Lohmann's universal QIK--VRT
correspondence thesis as an explicitly attributed scope claim, distinct from
the finite theorem, its physical instances, independent empirical
confirmation, and scientific consensus.
\end{abstract}

\noindent\textbf{Keywords:} QIK--VRT; observer-relative retrocausality;
local change time; Eigenzeit; provenance; distributed systems; information
direction; delayed choice; quantum eraser; recovery; responsibility.

\section{Purpose, status, and claim discipline}
\label{sec:purpose}

This is a successor manuscript, not an alteration of the earlier English
staging candidate.  The former candidate remains a byte-bound historical
intermediate state.  This version incorporates the subsequently clarified
point that \qikvrt{} \emph{Eigenzeit} is first an operational local
change-time concept, rather than a silent assertion that every use of the
word imports a relativistic metric.  The present source is local staging only
until an exact archive, destination metadata, author decision, platform
submission, and independent platform receipt have been recorded.

The paper makes several different types of statement.  They must not be
collapsed into one another:

\begin{itemize}
  \item A \emph{definition} specifies the relation called
  observer-relative retrocausality in this paper.
  \item A \emph{formal result} follows from explicitly declared
  assumptions.
  \item An \emph{executable witness} checks one finite, fully specified
  instance; it is not a proof-assistant kernel receipt or a survey of nature.
  \item A \emph{physical reference bridge} says what must be bound in a
  physical apparatus for the formal relation to be instantiated there.
  \item An \emph{empirical bridge} cites experimentally observed structures
  relevant to later context and earlier records.
  \item The \emph{universal correspondence thesis} is Ingolf Lohmann's
  attributed QIK--VRT claim about reality in its stated model scope.
\end{itemize}

Preserving these distinctions does not make the result weaker.  It makes its
scope inspectable and prevents a formal proof, a program run, a publication,
or an authority claim from being misdescribed as a different kind of
evidence.

\section{Difference, relation, and local change time}
\label{sec:local-time}

The QIK--VRT starting point is a relation of dependency:
\[
\begin{gathered}
\text{difference}\longrightarrow\text{information}
\longrightarrow\text{relation}\longrightarrow\text{causality}
\longrightarrow\text{spacetime}\\
\longrightarrow\text{appearance}\longrightarrow\text{life}
\longrightarrow\text{cognition}\longrightarrow\text{responsibility}
\longrightarrow\text{future}.
\end{gathered}
\]
In this vocabulary, information becomes relevant when a difference changes
possible connections, effects, or admissible descriptions.  This is the
operational reading of \emph{information efficacy}.  It is an ontological and
interpretative orientation of the research programme, not a claim that an
isolated slogan supersedes experimental method.

For the formal result, use a distributed event structure
\[
  H=(E,\host),
\]
where $E$ is a set of events and $\host$ is a strict partial order.  The
relation $e\host f$ means that $e$ precedes $f$ in the declared host or
effect-relevant order.  Irreflexivity and transitivity exclude a cycle in
$H$.  The structure is analogous to the causal partial orders used in
distributed systems \cite{Lamport1978}.

\begin{definition}[Observer and operational Eigenzeit]
\label{def:eigenzeit}
An observer or receiver $R$ is a chain $E_R\subseteq E$ of local events.  Its
local change time, operationally called \emph{Eigenzeit} in \qikvrt{}, is a
map
\[
  \tau_R:E_R\longrightarrow\mathbb{R}
\]
that is strictly order preserving on the chain:
\[
  e\host f\quad\Longrightarrow\quad \tau_R(e)<\tau_R(f)
  \qquad (e,f\in E_R).
\]
\end{definition}

An effective trigger, accepted receipt, provenance-bound record, or newly
available evidence can advance this local parameter.  A repository is an
ordinary example: the sequence of its effective, locally bound state changes
defines the relevant local order.  The definition does not require a global
clock or a spacetime metric.

For a physical receiver carried along a future-directed timelike worldline,
$\tau_R$ can additionally be calibrated to relativistic proper time.  In an
inertial frame, for $|v|<c$,
\[
  d\tau=dt\sqrt{1-v^2/c^2}>0.
\]
This metric fact is an additional physical binding, not a premise smuggled
into the operational definition.  Proper time on such a worldline does not
reverse sign \cite{Einstein1905,Minkowski1909}.  The theorem below needs the
shared monotonicity property, not a particular metric.

\section{Source order, information-bearing records, and the definition}
\label{sec:definition}

Let $S\subseteq E$ be a declared source chain.  A source mark
\[
  \theta:S\longrightarrow\mathbb{R}
\]
is strictly order preserving on $S$.  A record $r$ comprises at least a
source event $\sigma(r)\in S$, a payload $p(r)$, and a provenance/integrity
witness $w(r)$ binding the record to its source event.  Write
\[
  \theta(r)=\theta(\sigma(r)).
\]
The reception event at $R$ is $\rho_R(r)\in E_R$.  Every admissible record
satisfies the individual forward-delivery condition
\begin{equation}
  \sigma(r)\host\rho_R(r).
  \label{eq:forward}
\end{equation}

The binding matters.  Arbitrary labels on unrelated messages do not establish
an order relation.  The source mark must refer to a real, declared source
event and preserve the source-chain order.  The record must also be
information-bearing.  If $\mathcal{K}_{k-1}$ is the set of histories
compatible with the receiver's evidence before acceptance and
$\mathcal{K}_k$ the set afterwards, the acceptance is information-bearing
when
\begin{equation}
  \mathcal{K}_k\subsetneq\mathcal{K}_{k-1}.
  \label{eq:information}
\end{equation}

There are now at least two distinct orders on accepted records:
\begin{align}
 r\receive r' &\Longleftrightarrow
   \tau_R(\rho_R(r))<\tau_R(\rho_R(r')),\\
 r\sourceorder r' &\Longleftrightarrow
   \theta(r)<\theta(r').
\end{align}

\begin{definition}[Observer-relative retrocausality]
\label{def:orr}
Observer-relative retrocausality occurs for receiver $R$ when authentic,
information-bearing records $r_1,r_2$ satisfy
\begin{equation}
  r_1\receive r_2\qquad\text{and}\qquad r_2\sourceorder r_1.
  \label{eq:inverse}
\end{equation}
When numerical marks are declared, the comparative information-direction
slope is
\begin{equation}
  v_I^R(r_1,r_2)=
  \frac{\theta(r_2)-\theta(r_1)}
       {\tau_R(\rho_R(r_2))-\tau_R(\rho_R(r_1))}<0.
  \label{eq:slope}
\end{equation}
\end{definition}

The quotient in \cref{eq:slope} is an ordinal sign witness when a numerical
scaling is declared.  It is neither a signal speed nor a claimed Lorentz
invariant physical quantity.  The negative direction is a relation between
the local order in which a receiver changes and the provenance-bound source
order to which incoming information refers.

\section{Existence theorem and its limits}
\label{sec:theorem}

\begin{theorem}[Negative information direction]
\label{thm:existence}
Let $s^-\host s^+$ be two source events with
$\theta(s^-)<\theta(s^+)$.  Let $r^-$ and $r^+$ be authentic,
information-bearing records bound respectively to $s^-$ and $s^+$.  Suppose
their receiver-local delivery order is inverted:
\begin{equation}
  \rho_R(r^+)\host\rho_R(r^-),
  \label{eq:overtaking}
\end{equation}
while each record separately satisfies \cref{eq:forward}.  Then $R$ has
observer-relative retrocausality.  More specifically,
\[
  \Delta\tau_R=
  \tau_R(\rho_R(r^-)) - \tau_R(\rho_R(r^+))>0,
  \qquad
  \Delta\theta=\theta(r^-)-\theta(r^+)<0,
\]
and consequently $v_I^R(r^+,r^-)<0$.
\end{theorem}

\begin{proof}
Equation \cref{eq:overtaking} and the strict monotonicity of $\tau_R$ give
$\Delta\tau_R>0$.  Because the source mark is order preserving and
$s^-\host s^+$, its values obey $\theta(s^-)<\theta(s^+)$, hence
$\Delta\theta<0$.  The quotient in \cref{eq:slope} therefore has a negative
numerator and positive denominator.  Equation \cref{eq:forward} remains true
for each record, so the inverse relation lies between the source and
reception orders, not within the path of an individual record.
\end{proof}

\begin{corollary}[Order-preserving reparameterization]
\label{cor:reparameterization}
Replacing $\tau_R$ and $\theta$ by strictly increasing reparameterizations
preserves the inverse ordinal relation.
\end{corollary}

\begin{proof}
Strictly increasing maps preserve strict inequalities in both orders.
\end{proof}

\begin{proposition}[No causal loop follows]
\label{prop:no-loop}
The assumptions of \cref{thm:existence} do not imply a backward delivery of
one record or a causal loop in $H$.
\end{proposition}

\begin{proof}
For each record, \cref{eq:forward} says
$\sigma(r)\host\rho_R(r)$.  A reverse delivery of that same record would
also require $\rho_R(r)\host\sigma(r)$, contradicting irreflexivity after
transitivity.  A strict partial order has no closed causal chain.
\end{proof}

Thus the formally proved statement is precise: an inverse orientation between
two bound orders can coexist with monotonically increasing local change time,
forward delivery of every individual record, and acyclic host causality.  It
is not the proposition that a delayed packet is a time machine.

\section{Finite executable witness and positive-latency realization}
\label{sec:witness}

Take source marks $\theta(s^-)=10$ and $\theta(s^+)=20$.  Receiver $R$
accepts the newer-source record $r^+$ at local change time $100$ and the
older-source record $r^-$ at time $101$.  Then
\[
  v_I^R(r^+,r^-)=\frac{10-20}{101-100}=-10<0.
\]
Begin with four possible two-bit histories.  Accepting $r^+$ fixes one bit and
reduces the compatible histories from four to two.  Accepting $r^-$ fixes the
other and reduces them from two to one.  Thus both records meet
\cref{eq:information}.

The accompanying network-independent checker,
\path{verify_observer_relative_retrocausality.py}, evaluates this
finite witness and produces a canonical JSON report.  It checks host-order
acyclicity, source-before-reception, strict local monotonicity, strict
information-fibre refinement, inverse source/reception order, a second
observer with positive orientation, a future-directed latency construction,
and cancellation of complementary conditional fringes in an unconditioned
marginal.  It is an executable finite witness; the declared mathematical
proof remains the proof in \cref{sec:theorem}.

The positive-latency construction is elementary.  Let the old source event
occur at coordinate time $t^-=0$ with delay $d^-=3$, and the new source event
at $t^+=1$ with delay $d^+=1$.  The arrival times are
\[
  T^-=t^-+d^-=3,\qquad T^+=t^++d^+=2.
\]
Both delays are positive, but the later source event arrives first.  With
propagation at or below $c$, put $d^\pm=L^\pm/c$.  The sufficient overtaking
condition is
\begin{equation}
  \frac{L^- - L^+}{c}>t^+-t^-.
  \label{eq:latency}
\end{equation}
Unequal optical paths, network routes, queues, storage, or relays can
realize this construction.  No path is past-directed.

The same source records can also be received in source order by another
receiver $R'$.  That receiver has positive information direction.  The two
statements do not conflict because reception orders are local projections of
one distributed event structure; the source binding is held fixed.

A source observer's local-present event may be assigned to another observer's
coordinate future only when the compared events are spacelike separated.  This
coordinate assignment is not a causal-future relation: no record bound to that
event is available until a future-directed causal path reaches the receiver,
hence after its source emission.

\section{Physical bridge, delayed context, and no signalling}
\label{sec:physical}

The formal result is a physically instantiated operational relation only
when its terms are reference-bound.  A minimal bridge is:

\begin{center}
\begin{tabularx}{\textwidth}{>{\bfseries}p{0.30\textwidth}X}
\toprule
Formal term & Physical reference \\
\midrule
Receiver $R$ & physical receiving process or a timelike observer worldline \\
$\tau_R$ & monotonic local change time; optionally calibrated metric proper
time for a bound physical worldline \\
Source event $s$ & emission, measurement, or sealed record creation \\
$\theta(s)$ & pre-registered, order-preserving source mark \\
$\host$ & future-directed causal reachability or physical execution order \\
Record and witness & measured or stored record plus provenance/integrity binding \\
$\rho_R(r)$ & local receipt event bound to the receiver \\
\bottomrule
\end{tabularx}
\end{center}

If a physical trace satisfies this bridge and the inequalities of
\cref{thm:existence}, the defined relation is physically instantiated in that
trace.  This is not a claim that every physical system has the same structure
or that a particular apparatus proves a complete ontology.

Delayed-choice and quantum-eraser experiments give a complementary physical
pattern.  In an idealized eraser analysis, the conditional distributions can
have the form
\begin{align}
  P(x\mid j=1)&=\frac{1+\cos\phi(x)}{2},\\
  P(x\mid j=2)&=\frac{1-\cos\phi(x)}{2}.
\end{align}
At equal weights, their local unconditioned marginal is
\[
  P(x)=\tfrac12P(x\mid j=1)+\tfrac12P(x\mid j=2)=\tfrac12.
\]
The later partner/context record makes a correlation class of an earlier
signal record available; it does not rewrite the earlier signal record or
provide a selectable backward signal \cite{ScullyDruehl1982,Kim2000,Jordan1983}.
Related delayed-choice and causally disconnected-choice experiments maintain
the same distinction between joint contextual structure and controllable
communication \cite{Jacques2007,Ma2012,Ma2013}.

In this operational vocabulary, the no-signalling boundary is
\[
  C_{\leftarrow}=0,
\]
where $C_{\leftarrow}$ denotes the capacity of a controllable operational
channel into the receiver's causal past.  The boundary does not deny the
objective relevance of the joint correlation; it denies treating the
correlation as a usable backward message.

\section{Observation fibres, recovery, and responsibility}
\label{sec:recovery}

One local observation is not the whole event history.  Given an observation
operator $O$, the observation fibre is
\[
  K(o)=\{H\mid O(H)=o\}.
\]
When multiple histories are compatible with the same observation, no one of
them may be asserted as the uniquely actual history without further evidence.
Only claims that are invariant across the admissible fibre have the relevant
evidential status.  For a declared canonicalization $C$, recovery equivalence
has the form
\[
  H_1\sim_R H_2\quad\Longleftrightarrow\quad C(H_1)=C(H_2).
\]
Recovery therefore preserves what is responsibly reconstructible; it does
not license inventing a preferred past.

Later evidence can reduce the admissible fibre without overwriting earlier
records.  This is the same general distinction as in the delayed-context
case: the present grows by acquiring another effective relation, while the
stored earlier entry remains the entry it was.  The public QIK--VRT
declaration expresses the corresponding normative orientation as
\emph{enlarge the set of the present}: take more affected people,
consequences, evidence, and opportunities for correction into the current
field of responsibility.  This is a normative position, not a physical law.

The human stakes are concrete.  Incomplete evidence is often converted into
certainty, and unequal access to information is often converted into power.
Knowledge of relationships and expectations can be used for care, repair, and
liberation, or for manipulation, dependency, and domination.  A person must
not be reduced to material, risk, file, target group, enemy image, or a means
for an imposed future.  The normative QIK--VRT position is that responsibility
means preserving future agency.  It neither follows as a compulsory moral code
from the theorem nor grants anyone authority over another person.

\section{Correspondence thesis and explicit boundaries}
\label{sec:scope}

Ingolf Lohmann asserts the QIK--VRT correspondence thesis: the universe is a
distributed relational system and QIK--VRT describes reality within its
claimed model scope \cite{LohmannRoundTrip2026}.  The thesis is stated here
directly, neither hidden nor presented as an automatic consequence of the
finite theorem.  The theorem, a physical apparatus instance, the cited
empirical bridge, independent empirical confirmation, and scientific
consensus are distinct evidence states.

For a claimed instance of observer-relative retrocausality, the claim fails
if a source mark is unauthenticated or non-order-preserving, the two receipts
do not lie on the same monotone receiver chain, a purported record does not
refine admissible histories, the provenance binding can be silently altered,
or the inverse order disappears under strictly increasing
reparameterization.  These are testable failure conditions, not rhetorical
escape clauses.

The following are not claimed by the definition or theorem:
\begin{itemize}
  \item backward-running metric proper time;
  \item reception of a record before its own source emission;
  \item overwriting of a past raw record;
  \item a host causal loop;
  \item superluminal or past-directed propagation; or
  \item a controllable message into an observer's own causal past.
\end{itemize}

\section{Reproducibility and historical continuity}
\label{sec:reproducibility}

The public historical QIK--VRT corpus is preserved at Zenodo under
\zenodo.  The current German synthesis, its claim matrices, the finite
checker, and the canonical JSON witness are maintained in the QIK--VRT
repository.  Earlier English and German materials remain historical
intermediate states.  They are not silently overwritten or retroactively
given the scope wording of the present synthesis.

The local arXiv successor package contains the TeX source, a build receipt,
an exact manifest, a source archive, and an author-decision draft binding the
archive and destination metadata.  Until an arXiv account submission and a
returned arXiv receipt are verified, it does not establish an arXiv identifier,
arXiv acceptance, or any other external effect.

\section{Conclusion}

The operational result is compact:
\[
  \boxed{\Delta\tau_R>0\quad\text{and}\quad\Delta\theta<0.}
\]
The receiver changes forward in its local operational Eigenzeit while the
provenance-bound source references of newly received information can run in
the opposite direction.  This is observer-relative retrocausality as defined
here.  The relation is formally proved under disclosed assumptions, witnessed
by a finite executable checker, constructively realizable with forward
positive delays, and connected to physical experiments on later context and
earlier records.  It requires neither a reversed clock nor a message into the
past.

Its wider human consequence is not that a theorem decides morality.  It is
that reality, evidence, and future agency must not be severed from one
another.  We should distinguish, bind, preserve, correct, and take
responsibility for the relations through which our present becomes another
person's past or future.

\section*{Author contribution and disclosure}

Ingolf Lohmann originated the QIK--VRT research programme, the operational
terminology, the correspondence thesis, and the public normative declaration.
Artificial-intelligence assistance was used for English drafting, formal
organization, consistency checking, source-package construction, and layout.
The paper retains the distinction between author-attributed positions, formal
results, source-bound empirical references, normative statements, and open
scientific questions.

\begin{thebibliography}{99}

\bibitem{Einstein1905}
A. Einstein, Zur Elektrodynamik bewegter K\"orper, \emph{Annalen der
Physik} 322, 891--921 (1905).
\href{https://doi.org/10.1002/andp.19053221004}{doi:10.1002/andp.19053221004}.

\bibitem{Minkowski1909}
H. Minkowski, Raum und Zeit, \emph{Physikalische Zeitschrift} 10,
104--111 (1909).

\bibitem{Lamport1978}
L. Lamport, Time, Clocks, and the Ordering of Events in a Distributed
System, \emph{Communications of the ACM} 21, 558--565 (1978).
\href{https://doi.org/10.1145/359545.359563}{doi:10.1145/359545.359563}.

\bibitem{ScullyDruehl1982}
M. O. Scully and K. Dr\"uhl, Quantum eraser: A proposed photon correlation
experiment concerning observation and delayed choice in quantum mechanics,
\emph{Physical Review A} 25, 2208--2213 (1982).
\href{https://doi.org/10.1103/PhysRevA.25.2208}{doi:10.1103/PhysRevA.25.2208}.

\bibitem{Kim2000}
Y.-H. Kim, R. Yu, S. P. Kulik, Y. H. Shih, and M. O. Scully, Delayed
Choice Quantum Eraser, \emph{Physical Review Letters} 84, 1--5 (2000).
\href{https://doi.org/10.1103/PhysRevLett.84.1}{doi:10.1103/PhysRevLett.84.1}.

\bibitem{Jacques2007}
V. Jacques et al., Experimental Realization of Wheeler's Delayed-Choice
Gedanken Experiment, \emph{Science} 315, 966--968 (2007).
\href{https://doi.org/10.1126/science.1136303}{doi:10.1126/science.1136303}.

\bibitem{Ma2012}
X.-S. Ma et al., Experimental delayed-choice entanglement swapping,
\emph{Nature Physics} 8, 479--484 (2012).
\href{https://doi.org/10.1038/nphys2294}{doi:10.1038/nphys2294}.

\bibitem{Ma2013}
X.-S. Ma et al., Quantum erasure with causally disconnected choice,
\emph{Proceedings of the National Academy of Sciences} 110, 1221--1226
(2013). \href{https://doi.org/10.1073/pnas.1213201110}{doi:10.1073/pnas.1213201110}.

\bibitem{Jordan1983}
T. F. Jordan, Quantum correlations do not transmit signals,
\emph{Physics Letters A} 94, 264 (1983).
\href{https://doi.org/10.1016/0375-9601(83)90713-2}{doi:10.1016/0375-9601(83)90713-2}.

\bibitem{LohmannRoundTrip2026}
I. Lohmann, Von Softwarearchitektur zur Weltformel -- DAS UNIVERSUM ALS
ROUND TRIP, version 1.0.0-candidate.1, Zenodo (2026).
\href{https://doi.org/10.5281/zenodo.21888130}{doi:10.5281/zenodo.21888130}.

\end{thebibliography}

\end{document}
